% ============================================================
%  ANSWER KEY — Practice 8: Recurrence Relations
% ============================================================

\section*{Վարժություն 8. Անդրադարձ առնչություններ}

\subsubsection*{Խնդիր 1}
$a_1=3$, $a_2=9$, $a_{n+2}=a_{n+1}-a_n$:

Հաջորդականությունը պարբերական է 6 պարբերությամբ.
$3,\;9,\;6,\;-3,\;-9,\;-6,\;3,\;9,\;\ldots$

Քանի որ $200 = 33\times 6 + 2$, ունենք $a_{200}=a_2$:

$\boxed{a_{200}=9}$

\subsubsection*{Խնդիր 2}
$a_n=5a_{n-1}-6a_{n-2}$, $a_1=2$, $a_2=10$:

Բնութագրիչ հավասարում. $r^2-5r+6=0 \Rightarrow r=2,3$:
Ընդհանուր լուծում. $a_n=\alpha\cdot 2^n+\beta\cdot 3^n$:
Սկզբնական պայմաններից. $\alpha=-2$, $\beta=2$:

$\boxed{a_n = 2(3^n - 2^n)}$

\subsubsection*{Խնդիր 3}
$a_n=4a_{n-1}-4a_{n-2}$, $a_0=2$, $a_1=6$:

Բնութագրիչ հավասարում. $r^2-4r+4=0\Rightarrow r=2$ (կրկնակի արմատ):
Ընդհանուր լուծում. $a_n=(\alpha+\beta n)\cdot 2^n$:
$a_0=\alpha=2$, $a_1=(\alpha+\beta)\cdot 2 = 2(2+\beta)=6 \Rightarrow \beta=1$:

$\boxed{a_n = (2+n)\cdot 2^n}$

\subsubsection*{Խնդիր 4}
$a_n = 2\sqrt{2}\,a_{n-1}-4a_{n-2}$, $a_0=1$, $a_1=2$:

Բնութագրիչ հավասարում. $r^2-2\sqrt{2}\,r+4=0$:
$r = \sqrt{2}(1\pm i)$: Բևեռային տեսքով. $\rho=2$, $\theta=\pi/4$:

$a_0 = \alpha = 1$: $a_1 = \sqrt{2}(1+\beta)=2 \Rightarrow \beta = \sqrt{2}-1$:

$\boxed{a_n = 2^n\!\left(\cos\frac{n\pi}{4} + (\sqrt{2}-1)\sin\frac{n\pi}{4}\right)}$

\subsubsection*{Խնդիր 5}
$a_n=2a_{n-1}+a_{n-2}$, $a_0=0$, $a_1=1$:

Բնութագրիչ հավասարում. $r^2-2r-1=0 \Rightarrow r=1\pm\sqrt{2}$:
$a_0=\alpha+\beta=0$, $a_1=(1+\sqrt{2})\alpha+(1-\sqrt{2})\beta = 2\sqrt{2}\,\alpha=1$, ուստի $\alpha=\frac{1}{2\sqrt{2}}$:

$\boxed{a_n = \frac{(1+\sqrt{2})^n - (1-\sqrt{2})^n}{2\sqrt{2}}}$

\subsubsection*{Խնդիր 6}
10 երկարությամբ երկուական տողեր, որոնք պարունակում են (առնվազն) երկու հաջորդական զրոներ:

Հաշվել լրացումը. դիցուք $b_n$ = $n$ երկարությամբ տողեր առանց ``00''-ի:
$b_n = b_{n-1}+b_{n-2}$ (Ֆիբոնաչիի նման), $b_1=2$, $b_2=3$:
$b_{10}=144$: Ընդհանուր տողեր՝ $2^{10}=1024$:

$\boxed{1024-144=880}$

\subsubsection*{Խնդիր 7}
$n$ երկարությամբ բառեր $\{a,b,c\}$ բազմությունից այնպես, որ $a$-ն և $b$-ն երբեք հարևան չլինեն:

Դիցուք $A_n, B_n, C_n$ = բառեր, որոնք վերջանում են $a, b, c$-ով և չունեն հարևան $a,b$:
$A_n = A_{n-1}+C_{n-1}$, $B_n = B_{n-1}+C_{n-1}$, $C_n = A_{n-1}+B_{n-1}+C_{n-1}$:
$A_1=B_1=C_1=1$, ուստի $f(1)=3$:

Ընդհանուր քանակը $f(n) = A_n+B_n+C_n$ բավարարում է $f(n)=2f(n-1)+f(n-2)$ առնչությանը, որտեղ $f(1)=3, f(2)=7$:

\[
\boxed{f(n) = \frac{(1+\sqrt{2})^{n+1} + (1-\sqrt{2})^{n+1}}{2}}
\]

(Կոնկրետ երկարությունների համար. $f(10)=8119$:)